% Options for packages loaded elsewhere
\PassOptionsToPackage{unicode}{hyperref}
\PassOptionsToPackage{hyphens}{url}
%
\documentclass[
  ignorenonframetext,
  aspectratio=169,
]{beamer}
\usepackage{pgfpages}
\setbeamertemplate{caption}[numbered]
\setbeamertemplate{caption label separator}{: }
\setbeamercolor{caption name}{fg=normal text.fg}
\beamertemplatenavigationsymbolsempty
% Prevent slide breaks in the middle of a paragraph
\widowpenalties 1 10000
\raggedbottom
\setbeamertemplate{part page}{
  \centering
  \begin{beamercolorbox}[sep=16pt,center]{part title}
    \usebeamerfont{part title}\insertpart\par
  \end{beamercolorbox}
}
\setbeamertemplate{section page}{
  \centering
  \begin{beamercolorbox}[sep=12pt,center]{part title}
    \usebeamerfont{section title}\insertsection\par
  \end{beamercolorbox}
}
\setbeamertemplate{subsection page}{
  \centering
  \begin{beamercolorbox}[sep=8pt,center]{part title}
    \usebeamerfont{subsection title}\insertsubsection\par
  \end{beamercolorbox}
}
\AtBeginPart{
  \frame{\partpage}
}
\AtBeginSection{
  \ifbibliography
  \else
    \frame{\sectionpage}
  \fi
}
\AtBeginSubsection{
  \frame{\subsectionpage}
}

\usepackage{amsmath,amssymb}
\usepackage{iftex}
\ifPDFTeX
  \usepackage[T1]{fontenc}
  \usepackage[utf8]{inputenc}
  \usepackage{textcomp} % provide euro and other symbols
\else % if luatex or xetex
  \usepackage{unicode-math}
  \defaultfontfeatures{Scale=MatchLowercase}
  \defaultfontfeatures[\rmfamily]{Ligatures=TeX,Scale=1}
\fi
\usepackage{lmodern}
\usetheme[]{default}
\ifPDFTeX\else  
    % xetex/luatex font selection
\fi
% Use upquote if available, for straight quotes in verbatim environments
\IfFileExists{upquote.sty}{\usepackage{upquote}}{}
\IfFileExists{microtype.sty}{% use microtype if available
  \usepackage[]{microtype}
  \UseMicrotypeSet[protrusion]{basicmath} % disable protrusion for tt fonts
}{}
\makeatletter
\@ifundefined{KOMAClassName}{% if non-KOMA class
  \IfFileExists{parskip.sty}{%
    \usepackage{parskip}
  }{% else
    \setlength{\parindent}{0pt}
    \setlength{\parskip}{6pt plus 2pt minus 1pt}}
}{% if KOMA class
  \KOMAoptions{parskip=half}}
\makeatother
\usepackage{xcolor}
\newif\ifbibliography
\setlength{\emergencystretch}{3em} % prevent overfull lines
\setcounter{secnumdepth}{-\maxdimen} % remove section numbering

\usepackage{color}
\usepackage{fancyvrb}
\newcommand{\VerbBar}{|}
\newcommand{\VERB}{\Verb[commandchars=\\\{\}]}
\DefineVerbatimEnvironment{Highlighting}{Verbatim}{commandchars=\\\{\}}
% Add ',fontsize=\small' for more characters per line
\usepackage{framed}
\definecolor{shadecolor}{RGB}{241,243,245}
\newenvironment{Shaded}{\begin{snugshade}}{\end{snugshade}}
\newcommand{\AlertTok}[1]{\textcolor[rgb]{0.68,0.00,0.00}{#1}}
\newcommand{\AnnotationTok}[1]{\textcolor[rgb]{0.37,0.37,0.37}{#1}}
\newcommand{\AttributeTok}[1]{\textcolor[rgb]{0.40,0.45,0.13}{#1}}
\newcommand{\BaseNTok}[1]{\textcolor[rgb]{0.68,0.00,0.00}{#1}}
\newcommand{\BuiltInTok}[1]{\textcolor[rgb]{0.00,0.23,0.31}{#1}}
\newcommand{\CharTok}[1]{\textcolor[rgb]{0.13,0.47,0.30}{#1}}
\newcommand{\CommentTok}[1]{\textcolor[rgb]{0.37,0.37,0.37}{#1}}
\newcommand{\CommentVarTok}[1]{\textcolor[rgb]{0.37,0.37,0.37}{\textit{#1}}}
\newcommand{\ConstantTok}[1]{\textcolor[rgb]{0.56,0.35,0.01}{#1}}
\newcommand{\ControlFlowTok}[1]{\textcolor[rgb]{0.00,0.23,0.31}{#1}}
\newcommand{\DataTypeTok}[1]{\textcolor[rgb]{0.68,0.00,0.00}{#1}}
\newcommand{\DecValTok}[1]{\textcolor[rgb]{0.68,0.00,0.00}{#1}}
\newcommand{\DocumentationTok}[1]{\textcolor[rgb]{0.37,0.37,0.37}{\textit{#1}}}
\newcommand{\ErrorTok}[1]{\textcolor[rgb]{0.68,0.00,0.00}{#1}}
\newcommand{\ExtensionTok}[1]{\textcolor[rgb]{0.00,0.23,0.31}{#1}}
\newcommand{\FloatTok}[1]{\textcolor[rgb]{0.68,0.00,0.00}{#1}}
\newcommand{\FunctionTok}[1]{\textcolor[rgb]{0.28,0.35,0.67}{#1}}
\newcommand{\ImportTok}[1]{\textcolor[rgb]{0.00,0.46,0.62}{#1}}
\newcommand{\InformationTok}[1]{\textcolor[rgb]{0.37,0.37,0.37}{#1}}
\newcommand{\KeywordTok}[1]{\textcolor[rgb]{0.00,0.23,0.31}{#1}}
\newcommand{\NormalTok}[1]{\textcolor[rgb]{0.00,0.23,0.31}{#1}}
\newcommand{\OperatorTok}[1]{\textcolor[rgb]{0.37,0.37,0.37}{#1}}
\newcommand{\OtherTok}[1]{\textcolor[rgb]{0.00,0.23,0.31}{#1}}
\newcommand{\PreprocessorTok}[1]{\textcolor[rgb]{0.68,0.00,0.00}{#1}}
\newcommand{\RegionMarkerTok}[1]{\textcolor[rgb]{0.00,0.23,0.31}{#1}}
\newcommand{\SpecialCharTok}[1]{\textcolor[rgb]{0.37,0.37,0.37}{#1}}
\newcommand{\SpecialStringTok}[1]{\textcolor[rgb]{0.13,0.47,0.30}{#1}}
\newcommand{\StringTok}[1]{\textcolor[rgb]{0.13,0.47,0.30}{#1}}
\newcommand{\VariableTok}[1]{\textcolor[rgb]{0.07,0.07,0.07}{#1}}
\newcommand{\VerbatimStringTok}[1]{\textcolor[rgb]{0.13,0.47,0.30}{#1}}
\newcommand{\WarningTok}[1]{\textcolor[rgb]{0.37,0.37,0.37}{\textit{#1}}}

\providecommand{\tightlist}{%
  \setlength{\itemsep}{0pt}\setlength{\parskip}{0pt}}\usepackage{longtable,booktabs,array}
\usepackage{calc} % for calculating minipage widths
\usepackage{caption}
% Make caption package work with longtable
\makeatletter
\def\fnum@table{\tablename~\thetable}
\makeatother
\usepackage{graphicx}
\makeatletter
\def\maxwidth{\ifdim\Gin@nat@width>\linewidth\linewidth\else\Gin@nat@width\fi}
\def\maxheight{\ifdim\Gin@nat@height>\textheight\textheight\else\Gin@nat@height\fi}
\makeatother
% Scale images if necessary, so that they will not overflow the page
% margins by default, and it is still possible to overwrite the defaults
% using explicit options in \includegraphics[width, height, ...]{}
\setkeys{Gin}{width=\maxwidth,height=\maxheight,keepaspectratio}
% Set default figure placement to htbp
\makeatletter
\def\fps@figure{htbp}
\makeatother

% ============================================================
% StatLab Beamer Theme — Custom Header
% Mirrors the color scheme and layout of custom.scss (RevealJS)
% ============================================================

% ── COLORS (matches custom.scss variables) ───────────────────
\usepackage{xcolor}
\usepackage{pifont}

\definecolor{statprimary}{HTML}{2c3e50}    % dark gray-blue (body text)
\definecolor{statsecondary}{HTML}{2980b9}  % blue (structure, links)
\definecolor{statyale}{HTML}{00356b}       % Yale blue (footer)

% Box backgrounds
% Two-color system: Yale blue + statsecondary blue dominate.
% Orange (#FF6600) reserved for key warnings / key findings only.
\definecolor{stataccbg}{HTML}{e8f4f8}      % light cyan bg  — default block body
\definecolor{statorange}{HTML}{FF6600}     % orange         — alert / key finding header
\definecolor{statorangebg}{HTML}{FFF2E6}   % pale orange bg — alert body
\definecolor{statexbg}{HTML}{e6eef5}       % pale Yale-blue bg — example / Sam block body
\definecolor{statproof}{HTML}{95a5a6}      % gray           — proof (rarely used)
\definecolor{statproofbg}{HTML}{f9f9f9}    % light gray bg

% ── FONT: CMU Serif (LaTeX default) ─────────────────────────
% Switch Beamer from its default sans-serif to the classic serif.
\usefonttheme{serif}
% For XeLaTeX, explicitly select the OpenType CMU Serif face:
\setmainfont{CMU Serif}

% ── BEAMER COLORS ────────────────────────────────────────────

% Frame chrome
\setbeamercolor{frametitle}{bg=statprimary, fg=white}
\setbeamercolor{title}{fg=statprimary}
\setbeamercolor{subtitle}{fg=statsecondary}
\setbeamercolor{author}{fg=statprimary!70}
\setbeamercolor{date}{fg=statprimary!70}
\setbeamercolor{normal text}{fg=statprimary, bg=white}
\setbeamercolor{structure}{fg=statsecondary}

% Bullets and lists
\setbeamercolor{item}{fg=statsecondary}
\setbeamercolor{subitem}{fg=statsecondary!70}
\setbeamercolor{enumerate item}{fg=statsecondary}
\setbeamercolor{enumerate subitem}{fg=statsecondary!70}

% Bullet shape: filled circle (level 1), open circle (level 2), dash (level 3)
\setbeamertemplate{itemize item}{\small$\bullet$}
\setbeamertemplate{itemize subitem}{\small$\circ$}
\setbeamertemplate{itemize subsubitem}{\small\textendash}

% Default block — blue, mirrors .theorem / .lemma / .corollary
\setbeamercolor{block title}{fg=white, bg=statsecondary}
\setbeamercolor{block body}{bg=stataccbg, fg=statprimary}

% Alert block — orange (#FF6600), reserved for key warnings / findings
\setbeamercolor{block title alerted}{fg=white, bg=statorange}
\setbeamercolor{block body alerted}{bg=statorangebg, fg=statprimary}

% Example block — Yale blue, used for Sam callbacks and contextual examples
\setbeamercolor{block title example}{fg=white, bg=statyale}
\setbeamercolor{block body example}{bg=statexbg, fg=statprimary}

% ── FONT SIZES ───────────────────────────────────────────────
\setbeamerfont{title}{size=\LARGE, series=\bfseries}
\setbeamerfont{subtitle}{size=\large}
\setbeamerfont{frametitle}{size=\large, series=\bfseries}
\setbeamerfont{block title}{size=\normalsize, series=\bfseries}
\setbeamerfont{footline}{size=\scriptsize}

% ── FRAME TITLE ──────────────────────────────────────────────
\setbeamertemplate{navigation symbols}{}

\setbeamertemplate{frametitle}{%
  \nointerlineskip%
  \begin{beamercolorbox}[%
    wd=\paperwidth, ht=2.6ex, dp=0.9ex, leftskip=0.9em%
  ]{frametitle}%
    \usebeamerfont{frametitle}\insertframetitle%
  \end{beamercolorbox}%
}

% ── HEADLINE: presentation title, top-right (hidden on title slide) ──
\setbeamercolor{headline}{fg=statprimary!55, bg=white}
\setbeamertemplate{headline}{%
  \ifnum\insertframenumber>1%
    \begin{beamercolorbox}[wd=\paperwidth, ht=2.2ex, dp=0.5ex]{headline}%
      \hfill%
      \usebeamercolor[fg]{headline}{\fontsize{6.5}{8}\selectfont\inserttitle}%
      \hspace*{0.9em}%
    \end{beamercolorbox}%
  \else%
    \vspace{2.7ex}% reserve same height on title slide so frame title aligns
  \fi%
}

% ── FOOTER: branding + slide number (no logo) ────────────────
\setbeamertemplate{footline}{%
  {\color{statprimary!25}\hrule height 0.4pt}%
  \begin{beamercolorbox}[wd=\paperwidth, ht=2.4ex, dp=1.4ex]{footline}%
    \hspace*{0.9em}%
    \textcolor{statyale}{\fontfamily{qpl}\selectfont Yale Library $|$ StatLab}%
    \hfill%
    \textcolor{statprimary!50}{\insertframenumber\,/\,\inserttotalframenumber}%
    \hspace*{0.9em}%
  \end{beamercolorbox}%
}

% ── IMAGE CONSTRAINTS ────────────────────────────────────────
% Keep images within the content area so they never bleed into the footer.
% max-width = full text column; max-height = 65 % of the frame text height.
\usepackage{graphicx}
\setkeys{Gin}{width=\linewidth, height=0.65\textheight, keepaspectratio}

% ── TITLE PAGE ───────────────────────────────────────────────
\setbeamertemplate{title page}{%
  \vspace*{2.5em}%
  \begin{beamercolorbox}[sep=0pt, left]{title}%
    \usebeamerfont{title}\usebeamercolor[fg]{title}\inserttitle\par%
  \end{beamercolorbox}%
  \vspace*{0.5em}%
  \begin{beamercolorbox}[sep=0pt, left]{subtitle}%
    \usebeamerfont{subtitle}\usebeamercolor[fg]{subtitle}\insertsubtitle\par%
  \end{beamercolorbox}%
  \vspace*{1.8em}%
  \begin{beamercolorbox}[sep=0pt, left]{author}%
    \usebeamerfont{author}\usebeamercolor[fg]{author}\insertauthor\par%
  \end{beamercolorbox}%
  \vspace*{0.4em}%
  \begin{beamercolorbox}[sep=0pt, left]{date}%
    \usebeamerfont{date}\usebeamercolor[fg]{date}\insertdate\par%
  \end{beamercolorbox}%
}

% ── CUSTOM ENVIRONMENTS ──────────────────────────────────────
% Use these inside raw LaTeX blocks (```{=latex}) in your .qmd.
% Each mirrors the matching CSS class from the RevealJS theme.

% Definition — Yale blue (consistent with two-color palette)
\newenvironment{definitionblock}[1][Definition]{%
  \begingroup%
  \setbeamercolor{block title}{fg=white, bg=statyale}%
  \setbeamercolor{block body}{bg=statexbg, fg=statprimary}%
  \begin{block}{#1}%
}{%
  \end{block}%
  \endgroup%
}

% Lemma — blue, mirrors .lemma (same as default block)
\newenvironment{lemmablock}[1][Lemma]{%
  \begingroup%
  \begin{block}{#1}%
}{%
  \end{block}%
  \endgroup%
}

% Corollary — blue, mirrors .corollary (same as default block)
\newenvironment{corollaryblock}[1][Corollary]{%
  \begingroup%
  \begin{block}{#1}%
}{%
  \end{block}%
  \endgroup%
}

% Proof with QED square — gray, mirrors .proof
\newenvironment{proofblock}[1][Proof]{%
  \begingroup%
  \setbeamercolor{block title}{fg=white, bg=statproof}%
  \setbeamercolor{block body}{bg=statproofbg, fg=statprimary}%
  \begin{block}{\textit{#1}}%
  \itshape%
}{%
  \hfill$\square$%
  \end{block}%
  \endgroup%
}

% ── Workshop slide packages ───────────────────────────────────────────────
\usepackage{listings}
\usepackage{forest}
\usepackage{tikz}
\usetikzlibrary{arrows,arrows.meta,positioning,shapes.geometric,calc}
\usepackage{booktabs}
\usepackage{multicol}
\usepackage{array}
\usepackage{colortbl}   % \cellcolor / \rowcolor for highlighting table lookups

% ── Workshop color definitions ────────────────────────────────────────────
% Inline semantic colors (used in tables/code for before/after comparisons)
\definecolor{errorred}{RGB}{180,30,30}       % red  — bad examples, error states
\definecolor{highlightgreen}{RGB}{0,130,80}  % green — good examples, cached targets
\definecolor{commentgray}{RGB}{120,120,120}  % gray — code comments
\definecolor{warnorange}{RGB}{200,100,0}     % warm orange — string literals in code
\definecolor{lightblue}{RGB}{220,235,255}    % light blue — tikz fill

% ── Listings: tiny font, StatLab color accents ────────────────────────────
\lstset{
  basicstyle=\ttfamily\tiny,
  breaklines=true,
  frame=single,
  framerule=0.4pt,
  framesep=3pt,
  xleftmargin=8pt,
  xrightmargin=4pt,
  showstringspaces=false,
  commentstyle=\color{commentgray}\itshape,
  keywordstyle=\color{statsecondary}\bfseries,
  stringstyle=\color{warnorange},
  numbers=left,
  numberstyle=\tiny\color{gray!60},
  numbersep=5pt,
  tabsize=2,
  aboveskip=3pt,
  belowskip=3pt,
}

% ── Forest dirtree style ──────────────────────────────────────────────────
\forestset{
  dirtree/.style={
    for tree={
      font=\ttfamily\scriptsize,
      grow'=0,
      child anchor=west,
      parent anchor=south,
      anchor=west,
      calign=first,
      inner sep=1pt,
      outer sep=0pt,
      edge path={
        \noexpand\path [draw, \forestoption{edge}]
        (!u.south west) +(7.5pt,0) |- (.child anchor)
        \forestoption{edge label};
      },
      before typesetting nodes={
        if n=1{insert before={[,phantom]}}{}
      },
      fit=band,
      before computing xy={l=15pt},
    },
  }
}

% ── Quarto code block font size override ──────────────────────────────────
% Load fancyvrb explicitly so \DefineVerbatimEnvironment is always available,
% whether or not pandoc includes code blocks that trigger its auto-load.
\usepackage{fancyvrb}
\AtBeginDocument{%
  % Source code (syntax-highlighted input)
  \DefineVerbatimEnvironment{Highlighting}{Verbatim}{%
    commandchars=\\\{\},%
    fontsize=\small%
  }%
  % Printed results (plain verbatim) — matched to the input so a chunk
  % and its output read at one consistent size.
  \RecustomVerbatimEnvironment{verbatim}{Verbatim}{fontsize=\small}%
}

% ── Two-column shorthand colors for code comparisons ─────────────────────
\newcommand{\badlabel}{{\color{errorred}\textbf{$\boldsymbol{\times}$ Before}}}
\newcommand{\goodlabel}{{\color{highlightgreen}\textbf{\ding{51}~After}}}

% ── ORR Workshop rebuild macros (2025) ────────────────────────────────────

% Yale blue alias (for \color{yale})
\colorlet{yale}{statyale}

% Figure placeholder — reserves whitespace and flags human action needed
\newcommand{\figplaceholder}[1]{%
  \begin{center}%
  \fbox{\begin{minipage}{0.82\textwidth}%
  \centering\vspace{1.4cm}%
  {\large\color{gray} [FIGURE: #1]}\\[0.3cm]%
  {\small\color{gray}\itshape See [HUMAN INSTRUCTION] comment in source}%
  \vspace{1.4cm}%
  \end{minipage}}%
  \end{center}%
}

% Key finding callout (red alert block)
\newcommand{\findingbox}[1]{%
  \begin{alertblock}{Key Finding}%
  #1%
  \end{alertblock}%
}

% Title-less callout — example-block coloring, no title bar.
% Use where a block would otherwise invent a second heading under the frame title.
\newenvironment{notebox}{%
  \begin{center}%
  \begin{minipage}{0.95\textwidth}%
  \begin{beamercolorbox}[wd=\textwidth, sep=6pt]{block body example}%
}{%
  \end{beamercolorbox}%
  \end{minipage}%
  \end{center}%
}

% Inline pain/solution markers
\newcommand{\pain}[1]{{\color{errorred}\textbf{Problem:}} #1}
\newcommand{\solutiontag}[1]{{\color{yale}\textbf{Solution:}} #1}

% Parker Lab scenario callback — green example block
\newcommand{\samcallback}[1]{%
  \begin{exampleblock}{Parker Lab --- Sam's Problem}%
  \small #1%
  \end{exampleblock}%
}

% ── Named lstlisting styles ────────────────────────────────────────────────
% Override the global \tiny with \scriptsize for readability on slides.

\lstdefinestyle{Rcode}{%
  language=R,%
  basicstyle=\ttfamily\scriptsize,%
  backgroundcolor=\color{gray!8},%
  frame=single,%
  breaklines=true,%
  breakatwhitespace=false,%
  showstringspaces=false,%
  numbers=none,%
}

\lstdefinestyle{bashcode}{%
  language=bash,%
  basicstyle=\ttfamily\scriptsize,%
  backgroundcolor=\color{gray!8},%
  frame=single,%
  breaklines=true,%
  numbers=none,%
}

\lstdefinestyle{Pythoncode}{%
  language=Python,%
  basicstyle=\ttfamily\scriptsize,%
  backgroundcolor=\color{gray!8},%
  frame=single,%
  breaklines=true,%
  numbers=none,%
}

\lstdefinestyle{jsoncode}{%
  language={},%
  basicstyle=\ttfamily\scriptsize,%
  backgroundcolor=\color{gray!8},%
  frame=single,%
  breaklines=true,%
  numbers=none,%
  morestring=[b]",%
  stringstyle=\color{warnorange},%
}

\lstdefinestyle{yamlcode}{%
  language={},%
  basicstyle=\ttfamily\scriptsize,%
  backgroundcolor=\color{gray!8},%
  frame=single,%
  breaklines=true,%
  numbers=none,%
  morecomment=[l]{\#},%
  commentstyle=\color{commentgray}\itshape,%
}

\lstdefinestyle{textcode}{%
  language={},%
  basicstyle=\ttfamily\scriptsize,%
  backgroundcolor=\color{gray!8},%
  frame=none,%
  breaklines=true,%
  numbers=none,%
}
\makeatletter
\@ifpackageloaded{caption}{}{\usepackage{caption}}
\AtBeginDocument{%
\ifdefined\contentsname
  \renewcommand*\contentsname{Table of contents}
\else
  \newcommand\contentsname{Table of contents}
\fi
\ifdefined\listfigurename
  \renewcommand*\listfigurename{List of Figures}
\else
  \newcommand\listfigurename{List of Figures}
\fi
\ifdefined\listtablename
  \renewcommand*\listtablename{List of Tables}
\else
  \newcommand\listtablename{List of Tables}
\fi
\ifdefined\figurename
  \renewcommand*\figurename{Figure}
\else
  \newcommand\figurename{Figure}
\fi
\ifdefined\tablename
  \renewcommand*\tablename{Table}
\else
  \newcommand\tablename{Table}
\fi
}
\@ifpackageloaded{float}{}{\usepackage{float}}
\floatstyle{ruled}
\@ifundefined{c@chapter}{\newfloat{codelisting}{h}{lop}}{\newfloat{codelisting}{h}{lop}[chapter]}
\floatname{codelisting}{Listing}
\newcommand*\listoflistings{\listof{codelisting}{List of Listings}}
\makeatother
\makeatletter
\makeatother
\makeatletter
\@ifpackageloaded{caption}{}{\usepackage{caption}}
\@ifpackageloaded{subcaption}{}{\usepackage{subcaption}}
\makeatother
\ifLuaTeX
  \usepackage{selnolig}  % disable illegal ligatures
\fi
\usepackage{bookmark}

\IfFileExists{xurl.sty}{\usepackage{xurl}}{} % add URL line breaks if available
\urlstyle{same} % disable monospaced font for URLs
\hypersetup{
  pdftitle={Statistical Power},
  pdfauthor={Ted Ellsworth, PhD; Cole Brookson, MSc},
  hidelinks,
  pdfcreator={LaTeX via pandoc}}

\title{Statistical Power}
\author{Ted Ellsworth, PhD \and Cole Brookson, MSc}
\date{2026-10-15}

\begin{document}
\frame{\titlepage}

\begin{frame}{Free statistical consulting for Yale researchers}
\phantomsection\label{free-statistical-consulting-for-yale-researchers}
\vspace{2.8em}

\begin{center}

\large
Study design $\cdot$ analysis $\cdot$ visualization $\cdot$ writing.

\vspace{1.4em}

\normalsize
Walk-ins, appointments, and workshops --- free for all Yale affiliates.

\vspace{2.8em}

\textcolor{statsecondary}{%
  \large\textbf{library.yale.edu/statlab}%
}

\end{center}
\end{frame}

\begin{frame}{Introduction}
\phantomsection\label{introduction}
\vspace{1.4em}

\begin{center}
\large A committee, a reviewer, an IRB, or a funder asks: \textbf{how many subjects do you need?}
\end{center}

\begin{figure}

\centering{

\includegraphics[width=\textwidth,height=0.5\textheight]{./images/power-meme.png}

}

\caption{\label{fig-meme}All problems are solved with more power.}

\end{figure}%
\end{frame}

\begin{frame}{Schedule}
\phantomsection\label{schedule}
\vspace{0.3em}

Leaving this blank until we finalize
\end{frame}

\begin{frame}{Learning Objectives}
\phantomsection\label{learning-objectives}
\vspace{0.3em}

\small By the end of the session you should be able to:

\vspace{0.4em}

\begin{center}
\begin{enumerate}
  \item Describe the formal definition of statistical power, and relate it to other quantities of interest 
  \item Compute power by hand and with a simple online calculator
  \item Understand how to perform a "sweep" over a grid to identify a sample size 
  \item Generalize the formal definition of statistical power to an agnostic view of the "False negative" problem
  \item Calculate power from a closed form in the classic hypothesis testing framing 
  \item Understand and design an agnostic power calculation to quantify effect size recovery in an arbitrary model
\end{enumerate}
\end{center}
\end{frame}

\begin{frame}{False Negatives and False Positives (Type I and Type II
Errors)}
\phantomsection\label{false-negatives-and-false-positives-type-i-and-type-ii-errors}
\begin{columns}[T]
\begin{column}{0.6\textwidth}
\begin{itemize}
\item
  Most diagnostic statistics seek to identify two types of errors:

  \begin{itemize}
  \tightlist
  \item
    \emph{False Positive} (\(\alpha\), Type I) refer to the probability
    that you will claim an effect is \emph{present} when the
    \textbf{null hypothesis is true}.
  \item
    \emph{False Negatives} (\(\beta\), Type II) refer to the probability
    that you will claim an effect is \emph{absent} when the \textbf{null
    hypothesis is false}.
  \end{itemize}
\item
  Power is what is left over from a false negative error:
  \(\text{power} = 1 - \beta\).
\end{itemize}
\end{column}

\begin{column}{0.4\textwidth}
\begin{figure}

\centering{

\includegraphics[width=\textwidth,height=0.5\textheight]{./images/cohen.png}

}

\caption{\label{fig-cohen}Dr.~Jacob Cohen (1923-1998)}

\end{figure}%
\end{column}
\end{columns}

\begin{center}
Power is \textbf{the probability of rejecting the null hypothesis when it is in fact false and should be rejected.}
\end{center}
\end{frame}

\begin{frame}{Type S and Type M Errors}
\phantomsection\label{type-s-and-type-m-errors}
\begin{itemize}
\tightlist
\item
  Underpowered studies do more than miss real effects: when they
  \emph{do} find a significant effect, that estimate is likely to be
  wrong.
\item
  A significant estimate can be wrong in two ways:

  \begin{itemize}
  \tightlist
  \item
    \emph{Type S (sign)}: the estimate has the opposite sign of the true
    effect. At \(6\%\) power, \(24\%\) of significant estimates do.
  \item
    \emph{Type M (magnitude)}: the estimate overstates the true effect.
    At \(6\%\) power, significant estimates are \(9.5\) times too large
    on average.
  \end{itemize}
\end{itemize}

\vspace{0.4em}

\scriptsize Source: Gelman and Carlin (2014),
\textit{Beyond Power Calculations: Assessing Type S (Sign) and Type M (Magnitude) Errors}.
\end{frame}

\begin{frame}{What Power Looks Like}
\phantomsection\label{what-power-looks-like}
\begin{center}
\includegraphics{statistical-power-combined_files/figure-beamer/power-picture-1.png}
\end{center}

\vspace{-0.2em}

\begin{center}
\footnotesize \textcolor{errorred}{\textbf{Red}} is the study you would run if there were no effect; the shaded sliver past 1.96 is $\alpha$ --- the false positives you agreed to tolerate. \textcolor{statsecondary}{\textbf{Blue}} is the study you expect if the effect is real; \textbf{the shaded area past the same line is power}.
\end{center}
\end{frame}

\begin{frame}{Traditional Statistical Power \& Its Ingredients}
\phantomsection\label{traditional-statistical-power-its-ingredients}
\vskip0pt plus 1filll

So what do we need to actually calculate this? Well, the constituent
components to calculate \(1 - \beta\) are:

\begin{itemize}
\tightlist
\item
  \textbf{$\alpha$, the significance criterion}
\item
  \textbf{$N$, the sample size}
\item
  \textbf{$\delta$, the population effect size} \vskip0pt plus 1filll
\end{itemize}

\begin{center}
\textit{So,} you must be thinking, \textit{what's the formula?!}

\end{center}
\end{frame}

\begin{frame}{Traditional Statistical Power \& Its Ingredients}
\phantomsection\label{traditional-statistical-power-its-ingredients-1}
\begin{center}
\Large
There isn't a formula for the calculation of power. There are many!
\end{center}
\end{frame}

\begin{frame}{Core Notation}
\phantomsection\label{core-notation}
In order to estimate statistical power, we need four properties:

\vspace{0.5em}

\begin{center}
\small
\setlength{\tabcolsep}{7pt}
\renewcommand{\arraystretch}{1.35}
\begin{tabular}{c>{\raggedright\arraybackslash}p{3.1cm}>{\raggedright\arraybackslash}p{7.9cm}}
\toprule
\textbf{Symbol} & \textbf{Quality} & \textbf{What it is} \\
\midrule
$N$              & Sample size        & How many subjects you collect in total. \\
$\mu_t - \mu_c$  & Effect size        & The true difference between the treatment and control means: the \textit{signal}. \\
$\sigma$         & Standard deviation & How much outcomes vary from one another: the \textit{noise}. \\
$\alpha$         & Significance level & The Type I error rate you are willing to tolerate; it sets the evidence threshold. \\
\bottomrule
\end{tabular}
\end{center}

\vspace{0.4em}

\begin{center}
\small Note we will use $\mu_t - \mu_c$ throughout to denote effect size. Where space is tight --- in code, and in the appendix --- we abbreviate it $\delta$. 
\end{center}
\end{frame}

\begin{frame}{The Three Ingredients of Statistical Power}
\phantomsection\label{the-three-ingredients-of-statistical-power}
\vspace{0.5em}

\begin{center}
\small
\setlength{\tabcolsep}{7pt}
\renewcommand{\arraystretch}{1.35}
\begin{tabular}{>{\raggedright\arraybackslash}p{3.3cm}cc>{\raggedright\arraybackslash}p{6.2cm}}
\toprule
\textbf{Ingredient} & \textbf{Symbol} & \textbf{Control?} & \textbf{What it does} \\
\midrule
Strength of the treatment & $\mu_t - \mu_c$ & Partly & Dose, framing, and duration are yours to set; a stronger treatment is easier to detect. \\
Background noise & $\sigma$ & Partly & Noisier outcomes bury the signal, but better measurement and covariates quiet it. \\
Experimental design & $N$, $\alpha$ & \textbf{Yes} & Sample size, randomization scheme, and covariate adjustment. \\
\bottomrule
\end{tabular}
\end{center}
\end{frame}

\begin{frame}{Let's look at a case where we can find a closed form}
\phantomsection\label{lets-look-at-a-case-where-we-can-find-a-closed-form}
The formula depends entirely on what \emph{specific} test you want to
run. Here's an example. Let's make the following assumptions:

\begin{enumerate}
\tightlist
\item
  We have a completely randomized experiment, with two groups are
  selected to form a binary treatment, each group containing \(N \ 2\)
\item
  Assume the outcome variable is normally distributed, and the control
  group has mean \(\mu_c\) and the treatment group has \(\mu_t\)
\item
  The two groups have a shared standard deviation, \(\sigma\)
\end{enumerate}

We now say we want to perform a difference of means test, in the
two-tailed sense
\end{frame}

\begin{frame}{Let's look at a case where we can find a closed form}
\phantomsection\label{lets-look-at-a-case-where-we-can-find-a-closed-form-1}
Say we have the following hypothetical situation where we want to
compute power:

\begin{itemize}
\tightlist
\item
  A school is piloting a new \textbf{tutoring program}, and wants to
  know if it improves end-of-term test scores (out of 100)
\item
  Students are randomized into two groups, with one half getting
  tutoring
\item
  Untutored students average \(\mu_c = 72\)
\item
  The school will consider the program ``worth it'', if tutoring gets
  average scores to \(\mu_t = 75\)
\item
  Scores vary with \(\sigma = 10\)
\end{itemize}
\end{frame}

\begin{frame}{Let's look at a case where we can find a closed form}
\phantomsection\label{lets-look-at-a-case-where-we-can-find-a-closed-form-2}
Say we have the following hypothetical situation where we want to
compute power:

\begin{itemize}
\tightlist
\item
  A school is piloting a new \textbf{tutoring program}, and wants to
  know if it improves end-of-term test scores (out of 100)
\item
  Students are randomized into two groups, with one half getting
  tutoring
\item
  Untutored students average \(\mu_c = 72\)
\item
  The school will consider the program ``worth it'', if tutoring gets
  average scores to \(\mu_t = 75\)
\item
  Scores vary with \(\sigma = 10\)
\end{itemize}

\begin{center}
Luckily, in such a simple form, there's a nice asymptotic approximation for the power of the experiment
\end{center}
\end{frame}

\begin{frame}{Step 1: Do the Arithmetic}
\phantomsection\label{step-1-do-the-arithmetic}
\vspace{0.2em}

\begin{align*}
\text{power} &= \Phi\!\left( \frac{|\mu_t - \mu_c|\sqrt{N}}{2\sigma} \;-\; \Phi^{-1}\!\left(1-\frac{\alpha}{2}\right) \right)
   && \text{\small the formula} \\[0.25em]
&= \Phi\!\left( \frac{7\,\sqrt{164}}{2(20)} \;-\; \textcolor{statsecondary}{1.96} \right)
   && \text{\small plug in; \;\; $\Phi^{-1}(0.975) = 1.96$} \\[0.25em]
&= \Phi\!\left( \textcolor{statsecondary}{\mathbf{2.24}} \;-\; 1.96 \right) \;=\; \textcolor{statsecondary}{\Phi(\mathbf{0.28})}
   && \text{\small $7 \times 12.81 / 40 = 2.24$}
\end{align*}

\vspace{0.5em}

\begin{center}
\small
$\underbrace{\mathbf{2.24}}_{\text{what you expect}}$ \qquad
$\underbrace{\mathbf{1.96}}_{\text{what you need}}$ \qquad
$\underbrace{\mathbf{0.28}}_{\text{all the room you have}}$
\end{center}

\vspace{0.5em}

\begin{center}
\scriptsize A thin margin means an unreliable study. One lookup turns $\Phi(0.28)$ into a number. \\
(Where does $1.96$ come from? Running the $z$-table \textit{backwards} --- see the appendix.)
\end{center}
\end{frame}

\begin{frame}{Step 2: Read the Answer Off the Table}
\phantomsection\label{step-2-read-the-answer-off-the-table}
\begin{columns}[T]

\begin{column}{0.30\textwidth}
\vspace{0.6em}

\begin{center}
$\text{power} = \Phi(\mathbf{0.28})$
\end{center}

\vspace{0.5em}

\footnotesize
Split the $z$-score across the two margins of the table:

\vspace{0.4em}

\hspace{0.5em}$\mathbf{0.2}\;$ selects the \textbf{row}

\vspace{0.2em}

\hspace{0.5em}$\mathbf{0.08}$ selects the \textbf{column}

\vspace{0.5em}

The cell where they meet is the area to the \textit{left} of $z$, which is exactly the power.
\end{column}

\begin{column}{0.70\textwidth}
\begin{center}
\scriptsize
\setlength{\tabcolsep}{2.4pt}
\renewcommand{\arraystretch}{1.25}
\begin{tabular}{c|cccccccccc}
\toprule
$z$ & .00 & .01 & .02 & .03 & .04 & .05 & .06 & .07 & \cellcolor{stataccbg}\textbf{.08} & .09 \\
\midrule
0.0 & .5000 & .5040 & .5080 & .5120 & .5160 & .5199 & .5239 & .5279 & \cellcolor{stataccbg}.5319 & .5359 \\
0.1 & .5398 & .5438 & .5478 & .5517 & .5557 & .5596 & .5636 & .5675 & \cellcolor{stataccbg}.5714 & .5753 \\
\cellcolor{stataccbg}\textbf{0.2} & \cellcolor{stataccbg}.5793 & \cellcolor{stataccbg}.5832 & \cellcolor{stataccbg}.5871 & \cellcolor{stataccbg}.5910 & \cellcolor{stataccbg}.5948 & \cellcolor{stataccbg}.5987 & \cellcolor{stataccbg}.6026 & \cellcolor{stataccbg}.6064 & \cellcolor{statorange}\textbf{\color{white}.6103} & \cellcolor{stataccbg}.6141 \\
0.3 & .6179 & .6217 & .6255 & .6293 & .6331 & .6368 & .6406 & .6443 & \cellcolor{stataccbg}.6480 & .6517 \\
0.4 & .6554 & .6591 & .6628 & .6664 & .6700 & .6736 & .6772 & .6808 & \cellcolor{stataccbg}.6844 & .6879 \\
0.5 & .6915 & .6950 & .6985 & .7019 & .7054 & .7088 & .7123 & .7157 & \cellcolor{stataccbg}.7190 & .7224 \\
\bottomrule
\end{tabular}
\end{center}
\end{column}

\end{columns}

\vspace{0.5em}

\begin{center}
\large $\text{power} \;=\; \Phi(0.28) \;=\; \textcolor{statorange}{\mathbf{0.61}}$
\end{center}
\end{frame}

\begin{frame}{The Smallest Effect This Study Can See}
\phantomsection\label{the-smallest-effect-this-study-can-see}
\small We now know what this design does with a 7-point effect. The next
question is what it \textit{can} do: fix power at 80\% and solve the
same formula for the effect instead of for power.

\vspace{0.3em}

\begin{center}
$\displaystyle \text{MDE} \;=\; \underbrace{\left(\Phi^{-1}\!\left(1-\tfrac{\alpha}{2}\right) + \Phi^{-1}(0.80)\right)}_{1.96 \;+\; 0.84 \;=\; 2.80} \;\times\; \underbrace{\frac{2\sigma}{\sqrt{N}}}_{\text{standard error}}$
\end{center}

\vspace{0.3em}

\begin{center}
\scriptsize
\setlength{\tabcolsep}{9pt}
\renewcommand{\arraystretch}{1.08}
\begin{tabular}{ccc}
\toprule
$N$ & \textbf{Smallest detectable effect} & \textbf{In standard deviations} \\
\midrule
\textcolor{statorange}{\textbf{164}} & \textcolor{statorange}{\textbf{8.75 points}} & \textcolor{statorange}{\textbf{0.44$\sigma$}} \\
257  & 6.99 & 0.35$\sigma$ \\
400  & 5.60 & 0.28$\sigma$ \\
800  & 3.96 & 0.20$\sigma$ \\
1396 & 3.00 & 0.15$\sigma$ \\
\bottomrule
\end{tabular}
\end{center}

\vspace{0.3em}

\begin{center}
\small Pilot assumes $\mu_t - \mu_c = \mathbf{7}$. \qquad This design detects $\textcolor{statorange}{\mathbf{8.75}}$ and up.
\end{center}
\end{frame}

\begin{frame}{Check: Put 8.75 Back Through the Formula}
\phantomsection\label{check-put-8.75-back-through-the-formula}
\vspace{0.2em}

\begin{align*}
\text{power} &= \Phi\!\left( \frac{8.75\,\sqrt{164}}{2(20)} \;-\; 1.96 \right)
   && \text{\small the MDE in place of the pilot's effect} \\[0.25em]
&= \Phi\!\left( \textcolor{statsecondary}{\mathbf{2.80}} \;-\; 1.96 \right) \;=\; \Phi(\textcolor{statsecondary}{\mathbf{0.84}})
   && \text{\small $8.75 \times 12.81 / 40 = 2.80$} \\[0.25em]
&= \mathbf{0.80}
   && \text{\small row $0.8$, column $.04 \to .7995$}
\end{align*}

\vspace{0.4em}

\begin{center}
\footnotesize
\setlength{\tabcolsep}{9pt}
\renewcommand{\arraystretch}{1.15}
\begin{tabular}{>{\raggedright\arraybackslash}p{6.6cm}c}
\toprule
\textbf{Pin these three} & \textbf{Get} \\
\midrule
$N = 164$, \; $\sigma = 20$, \; effect $= 7$      & power $= 0.61$ \\
$N = 164$, \; $\sigma = 20$, \; power $= 0.80$    & effect $= 8.75$ \\
\bottomrule
\end{tabular}
\end{center}
\end{frame}

\begin{frame}{Your Turn: Is the Tutoring Program Worth Studying?}
\phantomsection\label{your-turn-is-the-tutoring-program-worth-studying}
\vspace{0.5em}

\begin{notebox}
\small A colleague is evaluating a new tutoring program. \textbf{400 students} will be randomly assigned, 200 to tutoring and 200 to business as usual, with final exam score as the outcome. Past cohorts averaged \textbf{72} points with a standard deviation of \textbf{10}, and the program is expected to add \textbf{3} points.
\end{notebox}

\vspace{0.7em}

\begin{center}
\large \textbf{What is the power of this design?}
\end{center}

\vspace{0.7em}

\begin{columns}[c]

\begin{column}{0.46\textwidth}
\begin{center}
\footnotesize
\setlength{\tabcolsep}{10pt}
\renewcommand{\arraystretch}{1.2}
\begin{tabular}{cc}
\toprule
\textbf{Symbol} & \textbf{Value} \\
\midrule
$N$      & 400  \\
$\mu_t$  & 75   \\
$\mu_c$  & 72   \\
$\sigma$ & 10   \\
$\alpha$ & 0.05, two-sided \\
\bottomrule
\end{tabular}
\end{center}
\end{column}

\begin{column}{0.50\textwidth}
\footnotesize
The formula and the $z$-table are on your \textbf{handout}.

\vspace{0.6em}

Work it in three moves:

\vspace{0.3em}

\hspace{0.4em}1. Compute what you \textbf{expect}

\vspace{0.2em}

\hspace{0.4em}2. Subtract what you \textbf{need}

\vspace{0.2em}

\hspace{0.4em}3. Look the difference up
\end{column}

\end{columns}
\end{frame}

\begin{frame}{Your Turn: The Solution}
\phantomsection\label{your-turn-the-solution}
\vspace{0.2em}

\begin{align*}
\text{power} &= \Phi\!\left( \frac{|75 - 72|\,\sqrt{400}}{2(10)} \;-\; \Phi^{-1}\!\left(1-\tfrac{0.05}{2}\right) \right)
   && \text{\small \textbf{Step 1:} plug in} \\[0.25em]
&= \Phi\!\left( \frac{3 \times 20}{20} \;-\; 1.96 \right)
   && \text{\small \textbf{Step 2:} $\sqrt{400} = 20$; $\Phi^{-1}(0.975) = 1.96$} \\[0.25em]
&= \Phi(3.00 - 1.96) \;=\; \textcolor{statsecondary}{\Phi(\mathbf{1.04})}
   && \text{\small \textbf{Step 3:} simplify} \\[0.25em]
&= \mathbf{0.85}
   && \text{\small \textbf{Step 4:} row $1.0$, column $.04 \to .8508$}
\end{align*}

\vspace{0.2em}
\end{frame}

\begin{frame}{Check Your Work: The EGAP Power Calculator}
\phantomsection\label{check-your-work-the-egap-power-calculator}
\vspace{0.3em}

\begin{center}
\small \url{https://egap.shinyapps.io/power-app/}
\end{center}

\vspace{0.6em}

\begin{center}
\small
\setlength{\tabcolsep}{7pt}
\renewcommand{\arraystretch}{1.25}
\begin{tabular}{c>{\raggedright\arraybackslash}p{6.4cm}c}
\toprule
\textbf{Our symbol} & \textbf{Field in the app} & \textbf{Enter} \\
\midrule
$\mu_t - \mu_c$ & Treatment Effect Size                  & 3    \\
$\sigma$        & Standard Deviation of Outcome Variable & 10   \\
$\alpha$        & Significance Level                     & 0.05 \\
                & Power Target                           & 0.8  \\
                & Maximum Number of Subjects             & 800  \\
\bottomrule
\end{tabular}
\end{center}

\vspace{0.5em}

\begin{center}
\small Leave \textit{Clustered Design} and \textit{Binary Dependent Variable} unchecked.
\end{center}
\end{frame}

\begin{frame}{Reading the Power Curve}
\phantomsection\label{reading-the-power-curve}
\begin{columns}[c]

\begin{column}{0.53\textwidth}
\includegraphics[width=\linewidth]{images/egap-power-app.png}
\end{column}

\begin{column}{0.45\textwidth}
\footnotesize
\setlength{\itemsep}{0.35em}
\begin{itemize}
\item The $x$-axis is \textbf{total sample size}; the $y$-axis is \textbf{power}.
\item The blue curve traces power as $N$ grows, holding the effect at 3 and $\sigma$ at 10.
\item The red dashed line marks the \textbf{80\% target}. Where the curve crosses it is the $N$ you need (here \textbf{349}).
\item Our $N = 400$ sits to the \textit{right} of that crossing. Reading up from 400 gives $\approx \mathbf{0.85}$.
\item The curve \textbf{flattens}: past $\sim$600 suggests diminishing returns at larger N.
\end{itemize}
\end{column}

\end{columns}
\end{frame}

\begin{frame}{Where to Situate Power Analysis in Your Research}
\phantomsection\label{where-to-situate-power-analysis-in-your-research}
\begin{itemize}
\tightlist
\item
  Power analysis is tricky because we need information about our data to
  derive power (e.g., \(\sigma\), \(\mu_t - \mu_c\)).
\item
  The problem is power calculations are most useful \emph{before} we've
  collected our data!
\item
  In order to do power analysis we need to do some educated guessing:

  \begin{itemize}
  \tightlist
  \item
    How big an effect should I anticipate?
  \item
    How many people can I realistically afford to treat?
  \item
    How many subjects will answer my survey?
  \end{itemize}
\end{itemize}
\end{frame}

\begin{frame}{What If We Don't Know the Effect?}
\phantomsection\label{what-if-we-dont-know-the-effect}
\vspace{0.2em}

\small Back to the survey experiment. The pilot fixed everything
\textit{except} for the anticipated effect size.

\vspace{0.5em}

\begin{columns}[T]

\begin{column}{0.40\textwidth}
\begin{center}
\footnotesize
\setlength{\tabcolsep}{6pt}
\renewcommand{\arraystretch}{1.25}
\begin{tabular}{cc}
\toprule
\multicolumn{2}{c}{\textbf{Fixed by the pilot}} \\
\midrule
$\mu_c$        & 50   \\
$\sigma$       & 20   \\
$\alpha$       & 0.05 \\
target power & 0.80 \\
\midrule
$\mu_t - \mu_c$ & \textcolor{statorange}{\textbf{?}} \\
$N$             & \textcolor{statorange}{\textbf{?}} \\
\bottomrule
\end{tabular}
\end{center}
\end{column}

\begin{column}{0.58\textwidth}
\begin{center}
\footnotesize
\setlength{\tabcolsep}{6pt}
\renewcommand{\arraystretch}{1.25}
\begin{tabular}{>{\raggedright\arraybackslash}p{5.6cm}c}
\toprule
\textbf{Where an effect could come from} & $\mu_t - \mu_c$ \\
\midrule
Our own pilot                          & 7 \\
A published study of a similar message & 5 \\
That published effect, discounted for publication bias & 3 \\
\bottomrule
\end{tabular}
\end{center}
\end{column}

\end{columns}

\vspace{0.5em}
\end{frame}

\begin{frame}[fragile]{The Formula, in Code}
\phantomsection\label{the-formula-in-code}
\vspace{0.6em}

\begin{center}
$\displaystyle \text{power} \;=\; \Phi\!\left( \frac{|\mu_t - \mu_c|\sqrt{N}}{2\sigma} \;-\; \Phi^{-1}\!\left(1-\frac{\alpha}{2}\right) \right)$
\end{center}

\vspace{0.9em}

\begin{Shaded}
\begin{Highlighting}[]
\NormalTok{power\_calc }\OtherTok{\textless{}{-}} \ControlFlowTok{function}\NormalTok{(mu\_t, mu\_c, sigma, N, }\AttributeTok{alpha =} \FloatTok{0.05}\NormalTok{) \{}
\NormalTok{    z }\OtherTok{\textless{}{-}} \FunctionTok{abs}\NormalTok{(mu\_t }\SpecialCharTok{{-}}\NormalTok{ mu\_c) }\SpecialCharTok{*} \FunctionTok{sqrt}\NormalTok{(N) }\SpecialCharTok{/}\NormalTok{ (}\DecValTok{2} \SpecialCharTok{*}\NormalTok{ sigma)}
\NormalTok{    cv }\OtherTok{\textless{}{-}} \FunctionTok{qnorm}\NormalTok{(}\DecValTok{1} \SpecialCharTok{{-}}\NormalTok{ alpha }\SpecialCharTok{/} \DecValTok{2}\NormalTok{)}
    \FunctionTok{pnorm}\NormalTok{(z }\SpecialCharTok{{-}}\NormalTok{ cv) }\SpecialCharTok{+} \FunctionTok{pnorm}\NormalTok{(}\SpecialCharTok{{-}}\NormalTok{z }\SpecialCharTok{{-}}\NormalTok{ cv) }\CommentTok{\# both rejection tails}
\NormalTok{\}}
\end{Highlighting}
\end{Shaded}

\vspace{0.9em}

\begin{center}
\large $\Phi \;=\;$ \texttt{pnorm} \qquad\qquad $\Phi^{-1} \;=\;$ \texttt{qnorm}
\end{center}
\end{frame}

\begin{frame}[fragile]{Power Across a Range of \(N\)}
\phantomsection\label{power-across-a-range-of-n}
\begin{Shaded}
\begin{Highlighting}[]
\NormalTok{grid }\OtherTok{\textless{}{-}} \FunctionTok{expand.grid}\NormalTok{(}\AttributeTok{N =} \FunctionTok{seq}\NormalTok{(}\DecValTok{60}\NormalTok{, }\DecValTok{1600}\NormalTok{, }\AttributeTok{by =} \DecValTok{10}\NormalTok{), }\AttributeTok{delta =} \FunctionTok{c}\NormalTok{(}\DecValTok{3}\NormalTok{, }\DecValTok{5}\NormalTok{, }\DecValTok{7}\NormalTok{))}

\NormalTok{grid}\SpecialCharTok{$}\NormalTok{power }\OtherTok{\textless{}{-}} \FunctionTok{power\_calc}\NormalTok{(}
    \AttributeTok{mu\_t =} \DecValTok{50} \SpecialCharTok{+}\NormalTok{ grid}\SpecialCharTok{$}\NormalTok{delta,}
    \AttributeTok{mu\_c =} \DecValTok{50}\NormalTok{,}
    \AttributeTok{sigma =} \DecValTok{20}\NormalTok{,}
    \AttributeTok{N =}\NormalTok{ grid}\SpecialCharTok{$}\NormalTok{N}
\NormalTok{)}

\FunctionTok{head}\NormalTok{(grid, }\DecValTok{3}\NormalTok{)}
\end{Highlighting}
\end{Shaded}

\begin{verbatim}
   N delta      power
1 60     3 0.08947303
2 70     3 0.09618740
3 80     3 0.10293349
\end{verbatim}

\vspace{0.4em}

\begin{center}
\small One formula, three assumptions, 465 rows. Now plot \texttt{power} against \texttt{N}.
\end{center}
\end{frame}

\begin{frame}{How Many Respondents Do We Need?}
\phantomsection\label{how-many-respondents-do-we-need}
\begin{center}
\includegraphics{statistical-power-combined_files/figure-beamer/sweep-plot-1.png}
\end{center}

\vspace{0.1em}
\end{frame}

\begin{frame}{Your Turn: Build the Sweep in R}
\phantomsection\label{your-turn-build-the-sweep-in-r}
\small Your department is considering a six-week \textbf{peer-mentoring
program} for first-year students, and can randomize students into it or
onto a waitlist. The outcome is end-of-year GPA. The registrar tells you
first-year GPA averages \textbf{3.05}, with a standard deviation of
\textbf{0.45}.

\vspace{0.5em}

\begin{columns}[T]
\begin{column}{0.44\textwidth}
\begin{center}
\footnotesize
\setlength{\tabcolsep}{6pt}
\renewcommand{\arraystretch}{1.25}
\begin{tabular}{cc}
\toprule
\multicolumn{2}{c}{\textbf{Given}} \\
\midrule
$\mu_c$  & 3.05 \\
$\sigma$ & 0.45 \\
$\alpha$ & 0.05 \\
\bottomrule
\end{tabular}
\end{center}
\end{column}

\begin{column}{0.54\textwidth}
\vspace{0.4em}
\footnotesize
\textbf{You decide:}

\vspace{0.3em}

\hspace{0.4em}\textbullet~Which GPA gains are plausible?

\vspace{0.2em}

\hspace{0.4em}\textbullet~How many students could you enroll?
\end{column}
\end{columns}

\vspace{0.4em}
\end{frame}

\begin{frame}[fragile]{Your Turn: One Solution}
\phantomsection\label{your-turn-one-solution}
\begin{Shaded}
\begin{Highlighting}[]
\NormalTok{grid }\OtherTok{\textless{}{-}} \FunctionTok{expand.grid}\NormalTok{(}\AttributeTok{N =} \FunctionTok{seq}\NormalTok{(}\DecValTok{50}\NormalTok{, }\DecValTok{1000}\NormalTok{, }\AttributeTok{by =} \DecValTok{10}\NormalTok{), }\AttributeTok{delta =} \FunctionTok{c}\NormalTok{(}\FloatTok{0.10}\NormalTok{, }\FloatTok{0.15}\NormalTok{, }\FloatTok{0.20}\NormalTok{))}
\NormalTok{grid}\SpecialCharTok{$}\NormalTok{power }\OtherTok{\textless{}{-}} \FunctionTok{power\_calc}\NormalTok{(}
    \AttributeTok{mu\_t =} \FloatTok{3.05} \SpecialCharTok{+}\NormalTok{ grid}\SpecialCharTok{$}\NormalTok{delta,}
    \AttributeTok{mu\_c =} \FloatTok{3.05}\NormalTok{,}
    \AttributeTok{sigma =} \FloatTok{0.45}\NormalTok{,}
    \AttributeTok{N =}\NormalTok{ grid}\SpecialCharTok{$}\NormalTok{N}
\NormalTok{)}

\ControlFlowTok{for}\NormalTok{ (d }\ControlFlowTok{in} \FunctionTok{unique}\NormalTok{(grid}\SpecialCharTok{$}\NormalTok{delta)) \{}
\NormalTok{    ok }\OtherTok{\textless{}{-}}\NormalTok{ grid}\SpecialCharTok{$}\NormalTok{N[grid}\SpecialCharTok{$}\NormalTok{delta }\SpecialCharTok{==}\NormalTok{ d }\SpecialCharTok{\&}\NormalTok{ grid}\SpecialCharTok{$}\NormalTok{power }\SpecialCharTok{\textgreater{}=} \FloatTok{0.80}\NormalTok{]}
    \FunctionTok{cat}\NormalTok{(}\StringTok{"effect"}\NormalTok{, d, }\StringTok{"GPA points {-}\textgreater{} N ="}\NormalTok{, }\FunctionTok{min}\NormalTok{(ok), }\StringTok{"}\SpecialCharTok{\textbackslash{}n}\StringTok{"}\NormalTok{)}
\NormalTok{\}}
\end{Highlighting}
\end{Shaded}

\begin{verbatim}
effect 0.1 GPA points -> N = 640 
effect 0.15 GPA points -> N = 290 
effect 0.2 GPA points -> N = 160 
\end{verbatim}

\footnotesize Your effects and \(N\) range will differ. The
\emph{structure} is the answer: one grid, one call, one threshold.
\end{frame}

\begin{frame}{But what about literally any other case?}
\phantomsection\label{but-what-about-literally-any-other-case}
There are lots of reasons why, even for simple problems, we don't want
to rely on such an approximation. For instance, what if you

\begin{enumerate}
\tightlist
\item
  don't assume a two-tailed test?
\item
  have clustering and want some sort of random effect?
\item
  are taking a Bayesian approach?
\item
  are using nonparametric statistics?
\end{enumerate}
\end{frame}

\begin{frame}{What should we do in those other cases?}
\phantomsection\label{what-should-we-do-in-those-other-cases}
\begin{itemize}
\tightlist
\item
  Even when power is not as clearly defined, there are still a whole
  range of things we still want to know about ``power''
\item
  Let's start relaxing our definition of power. We can, for now,
  consider it as:
\end{itemize}

\vskip50pt

\begin{center}
\textbf{The probability that we we \textcolor{red}{recover} a real effect, should one \textcolor{red}{exist}}
\end{center}
\end{frame}

\begin{frame}{What should we do in those other cases?}
\phantomsection\label{what-should-we-do-in-those-other-cases-1}
What would it mean then, to approach asking questions about power or
sample size in a more general sense where we may want to consider:

\vskip20pt

\begin{columns}[T]
\begin{column}{0.6\textwidth}
\begin{itemize}
\tightlist
\item
  Non-parametric approaches
\item
  Bayesian approaches
\item
  Or even just moderately more complicated study designs?
\end{itemize}

\vskip20pt

\begin{center}
\textbf{Well, we go gambling}
\end{center}
\end{column}

\begin{column}{0.4\textwidth}
\begin{figure}

\centering{

\includegraphics[width=\textwidth,height=0.5\textheight]{./images/casino.png}

}

\caption{\label{fig-casino}Oxford Casino Hotel}

\end{figure}%
\end{column}
\end{columns}
\end{frame}

\begin{frame}{An aside into the casinos at Monaco}
\phantomsection\label{an-aside-into-the-casinos-at-monaco}
\begin{columns}[T]
\begin{column}{0.6\textwidth}
\begin{itemize}
\tightlist
\item
  Monte Carlo refers to a broad class of methods wherein you use
  repeated random sampling to obtain some form of numerical result
\item
  It will be very useful for us when we cannot rely on a closed form
  solution (which, as a reminder, is in most cases!)
\item
  Essentially, we can use \textcolor{red}{randomness} to solve
  \textcolor{red}{deterministic} problems
\end{itemize}

\vskip20pt

\textit{For any problem with probabilistic interpretation, it can be solved with Monte Carlo}
\end{column}

\begin{column}{0.4\textwidth}
\begin{figure}

\centering{

\includegraphics[width=\textwidth,height=0.5\textheight]{./images/ulam.png}

}

\caption{\label{fig-casino}Primary developer of the MC method, Stanisław
Ulam, who was inspired by his uncle's gambling}

\end{figure}%
\end{column}
\end{columns}
\end{frame}

\begin{frame}{An aside into the casinos at Monaco}
\phantomsection\label{an-aside-into-the-casinos-at-monaco-1}
\begin{columns}[T]
\begin{column}{0.6\textwidth}
Let's imagine we want to know the answer to a gambling-related problem
with a closed-form answer:

\begin{center}
What's the probability you win on red at the roulette table three spins in a row?
\end{center}

\begin{itemize}
\tightlist
\item
  There are 18 \textcolor{red}{red} pockets out of 38, so the exact
  answer is
\end{itemize}

\[
P(\text{3 reds in a row}) = (18/38)^3 \approx 0.106
\]
\end{column}

\begin{column}{0.4\textwidth}
\begin{figure}

\centering{

\includegraphics[width=\textwidth,height=0.5\textheight]{./images/roulette.png}

}

\caption{\label{fig-casino}A reasonably standard roulette table}

\end{figure}%
\end{column}
\end{columns}
\end{frame}

\begin{frame}[fragile]{An aside to the casinos at Monaco}
\phantomsection\label{an-aside-to-the-casinos-at-monaco}
But we could also just\ldots{} go to the casino 100,000 times:

\begin{Shaded}
\begin{Highlighting}[]
\FunctionTok{set.seed}\NormalTok{(}\DecValTok{2026}\NormalTok{)}
\NormalTok{spins }\OtherTok{\textless{}{-}} \FunctionTok{matrix}\NormalTok{(}\FunctionTok{runif}\NormalTok{(}\DecValTok{3} \SpecialCharTok{*} \DecValTok{100000}\NormalTok{) }\SpecialCharTok{\textless{}} \DecValTok{18} \SpecialCharTok{/} \DecValTok{38}\NormalTok{, }\AttributeTok{ncol =} \DecValTok{3}\NormalTok{) }\CommentTok{\# each row is 3 spins}
\FunctionTok{mean}\NormalTok{(}\FunctionTok{rowSums}\NormalTok{(spins) }\SpecialCharTok{==} \DecValTok{3}\NormalTok{) }\CommentTok{\# proportion of rows where all three won}
\end{Highlighting}
\end{Shaded}

\begin{verbatim}
[1] 0.1059
\end{verbatim}

\begin{center}
If we can describe \textcolor{red}{how the data are generated}, we can estimate \textit{any} probability by running the process many times!
\end{center}
\end{frame}

\begin{frame}{An aside into the casinos at Monaco}
\phantomsection\label{an-aside-into-the-casinos-at-monaco-2}
\vskip0pt plus 1filll

\begin{center}
\Large
And whether its the more traditional, rigid form, or our more flexible modern form, power is \textit{just a probability.}
\end{center}

\vskip0pt plus 1filll
\end{frame}

\begin{frame}{So, how do we do a Monte Carlo simulation, anyways?}
\phantomsection\label{so-how-do-we-do-a-monte-carlo-simulation-anyways}
In the parametric context, a Monte Carlo power analysis looks like this:

\begin{enumerate}
\tightlist
\item
  \textbf{Generate} random data from some distribution with the
  properties that you expect your real data to have
\item
  \textbf{Analyze} this randomly generated data set with your
  statistical analysis, whatever that is
\item
  \textbf{Store} the value(s) of the analysis that you're interested in
\item
  \textbf{Repeat} the above steps many times
\item
  \textbf{Compute} power by looking at the proportion of times the
  statistic or value you're interested in was ``correct''
\end{enumerate}

\vskip20pt

\begin{center}
This can be arbitrarily complex or simple!
\end{center}
\end{frame}

\begin{frame}[fragile]{So, how do we do a Monte Carlo simulation,
anyways?}
\phantomsection\label{so-how-do-we-do-a-monte-carlo-simulation-anyways-1}
We can consider all of this to be done, for the sake of simplicity, by
three piecemeal functions that we will want to write

\begin{itemize}
\tightlist
\item
  \texttt{generate()}: make a fake data set that looks like your study
\item
  \texttt{analyze()}: run the analysis you'd \emph{actually} run
\item
  \texttt{success()}: decide whether this run ``recovered'' the effect
\end{itemize}

and one function that necessarily performs the actual simulation

\begin{Shaded}
\begin{Highlighting}[]
\NormalTok{simulate\_power }\OtherTok{\textless{}{-}} \ControlFlowTok{function}\NormalTok{(n\_sims, generate, analyze, success) \{}
\NormalTok{    hits }\OtherTok{\textless{}{-}} \FunctionTok{replicate}\NormalTok{(n\_sims, }\FunctionTok{success}\NormalTok{(}\FunctionTok{analyze}\NormalTok{(}\FunctionTok{generate}\NormalTok{())))}
    \FunctionTok{mean}\NormalTok{(hits) }\CommentTok{\# proportion of successful runs = power}
\NormalTok{\}}
\end{Highlighting}
\end{Shaded}
\end{frame}

\begin{frame}{So, how do we do a Monte Carlo simulation, anyways?}
\phantomsection\label{so-how-do-we-do-a-monte-carlo-simulation-anyways-2}
We need to make a statement on how we think the data are being
generated. For student \(i\):

\[
\text{post}_i = \underbrace{\mu_c}_{\text{baseline}} + \underbrace{\delta \, T_i}_{\text{tutoring effect}} + \underbrace{\varepsilon_i}_{\text{noise}}, \qquad \varepsilon_i \sim N(0, \sigma^2)
\]

\vskip10pt

\begin{itemize}
\tightlist
\item
  \(\mu_c = 72\): the average score \emph{without} tutoring
\item
  \(T_i\): 1 if student \(i\) was tutored, 0 if not, so tutored students
  get an ``effect'' of \(\delta = 3\)
\item
  \(\varepsilon_i\): our random noise with \(\sigma = 10\)
\end{itemize}

\vskip10pt
\end{frame}

\begin{frame}[fragile]{So, how do we do a Monte Carlo simulation,
anyways?}
\phantomsection\label{so-how-do-we-do-a-monte-carlo-simulation-anyways-3}
Let's start with our function that generates data based on our student
tutoring example from before.

\begin{Shaded}
\begin{Highlighting}[]
\NormalTok{gen\_basic }\OtherTok{\textless{}{-}} \ControlFlowTok{function}\NormalTok{(}\AttributeTok{n =} \DecValTok{350}\NormalTok{, }\AttributeTok{effect =} \DecValTok{3}\NormalTok{, }\AttributeTok{sd =} \DecValTok{10}\NormalTok{) \{}
\NormalTok{    treat }\OtherTok{\textless{}{-}} \FunctionTok{rep}\NormalTok{(}\DecValTok{0}\SpecialCharTok{:}\DecValTok{1}\NormalTok{, }\AttributeTok{each =}\NormalTok{ n }\SpecialCharTok{/} \DecValTok{2}\NormalTok{) }\CommentTok{\# half control, half tutored}
\NormalTok{    post }\OtherTok{\textless{}{-}} \DecValTok{72} \SpecialCharTok{+}\NormalTok{ effect }\SpecialCharTok{*}\NormalTok{ treat }\SpecialCharTok{+} \FunctionTok{rnorm}\NormalTok{(n, }\AttributeTok{mean =} \DecValTok{0}\NormalTok{, }\AttributeTok{sd =}\NormalTok{ sd)}
    \FunctionTok{data.frame}\NormalTok{(treat, post)}
\NormalTok{\}}

\FunctionTok{head}\NormalTok{(}\FunctionTok{gen\_basic}\NormalTok{())}
\end{Highlighting}
\end{Shaded}

\begin{verbatim}
  treat     post
1     0 60.03729
2     0 86.41172
3     0 82.52832
4     0 66.93240
5     0 74.34064
6     0 72.20875
\end{verbatim}
\end{frame}

\begin{frame}{So, how do we do a Monte Carlo simulation, anyways?}
\phantomsection\label{so-how-do-we-do-a-monte-carlo-simulation-anyways-4}
\includegraphics{statistical-power-combined_files/figure-beamer/unnamed-chunk-4-1.pdf}
\end{frame}

\begin{frame}[fragile]{So, how do we do a Monte Carlo simulation,
anyways?}
\phantomsection\label{so-how-do-we-do-a-monte-carlo-simulation-anyways-5}
\begin{Shaded}
\begin{Highlighting}[]
\NormalTok{analyze\_lm }\OtherTok{\textless{}{-}} \ControlFlowTok{function}\NormalTok{(dat) \{}
\NormalTok{  fit }\OtherTok{\textless{}{-}} \FunctionTok{lm}\NormalTok{(post }\SpecialCharTok{\textasciitilde{}}\NormalTok{ treat, }\AttributeTok{data =}\NormalTok{ dat) }\CommentTok{\# same as an equal{-}variance t{-}test}
  \FunctionTok{summary}\NormalTok{(fit)}\SpecialCharTok{$}\NormalTok{coefficients[}\StringTok{"treat"}\NormalTok{, ] }\CommentTok{\# estimate, SE, t, p for treat}
\NormalTok{\}}

\NormalTok{success\_sig }\OtherTok{\textless{}{-}} \ControlFlowTok{function}\NormalTok{(est) \{}
\NormalTok{  est[[}\StringTok{"Pr(\textgreater{}|t|)"}\NormalTok{]] }\SpecialCharTok{\textless{}} \FloatTok{0.05} \CommentTok{\# did we reject the null?}
\NormalTok{\}}

\FunctionTok{analyze\_lm}\NormalTok{(}\FunctionTok{gen\_basic}\NormalTok{())}
\end{Highlighting}
\end{Shaded}

\begin{verbatim}
  Estimate Std. Error    t value   Pr(>|t|) 
2.48850940 1.10064783 2.26094971 0.02437943 
\end{verbatim}
\end{frame}

\begin{frame}[fragile]{Does the casino agree with our nice equation??}
\phantomsection\label{does-the-casino-agree-with-our-nice-equation}
\begin{Shaded}
\begin{Highlighting}[]
\FunctionTok{set.seed}\NormalTok{(}\DecValTok{2026}\NormalTok{)}
\FunctionTok{simulate\_power}\NormalTok{(}
    \AttributeTok{n\_sims =} \DecValTok{2000}\NormalTok{,}
    \AttributeTok{generate =}\NormalTok{ gen\_basic,}
    \AttributeTok{analyze =}\NormalTok{ analyze\_lm,}
    \AttributeTok{success =}\NormalTok{ success\_sig}
\NormalTok{)}
\end{Highlighting}
\end{Shaded}

\begin{verbatim}
[1] 0.8105
\end{verbatim}

\vskip10pt

\begin{itemize}
\tightlist
\item
  Pretty close for only 2000 simulations!
\end{itemize}
\end{frame}

\begin{frame}{Does the casino agree with our nice equation??}
\phantomsection\label{does-the-casino-agree-with-our-nice-equation-1}
\begin{itemize}
\tightlist
\item
  It's worth noting that the simulation has its own error! With \(S\)
  simulations, the Monte Carlo standard error is
\end{itemize}

\[
\text{SE} = \sqrt{\frac{p(1-p)}{S}}
\]

\begin{center}
Want a more precise answer? \textbf{Go to the casino more times.}
\end{center}
\end{frame}

\begin{frame}{What if we make this more complicated, arbitrarily so?}
\phantomsection\label{what-if-we-make-this-more-complicated-arbitrarily-so}
I suspect that schools know a lot about their students \emph{before} the
program starts. Say we also have each student's \textbf{pre-test score}.
Our model for student \(i\) becomes:

\[
\text{pre}_i \sim N(70, \sigma^2)
\]

\[
\text{post}_i = \underbrace{\mu_c}_{\text{baseline}} + 
\underbrace{\rho\,(\text{pre}_i - 70)}_{\text{what we knew already}} + 
\underbrace{\delta \, T_i}_{\text{tutoring effect}} + 
\underbrace{\varepsilon_i}_{\text{noise}}, 
\qquad \varepsilon_i \sim N\big(0,\ \sigma^2(1-\rho^2)\big)
\]

\vskip05pt

\begin{itemize}
\tightlist
\item
  \(\rho = 0.6\) is the correlation between pre-test and end-of-term
  scores
\item
  Students who were above average before tend to stay above average
\item
  The total spread of scores is still \(\sigma = 10\)
\end{itemize}
\end{frame}

\begin{frame}[fragile]{What if we make this more complicated,
arbitrarily so?}
\phantomsection\label{what-if-we-make-this-more-complicated-arbitrarily-so-1}
\begin{itemize}
\tightlist
\item
  We only need to change \textbf{two} of our three pieces:
\end{itemize}

\begin{Shaded}
\begin{Highlighting}[]
\NormalTok{gen\_pretest }\OtherTok{\textless{}{-}} \ControlFlowTok{function}\NormalTok{(}\AttributeTok{n =} \DecValTok{350}\NormalTok{, }\AttributeTok{effect =} \DecValTok{3}\NormalTok{, }\AttributeTok{sd =} \DecValTok{10}\NormalTok{, }\AttributeTok{rho =} \FloatTok{0.6}\NormalTok{) \{}
\NormalTok{  treat }\OtherTok{\textless{}{-}} \FunctionTok{rep}\NormalTok{(}\DecValTok{0}\SpecialCharTok{:}\DecValTok{1}\NormalTok{, }\AttributeTok{each =}\NormalTok{ n }\SpecialCharTok{/} \DecValTok{2}\NormalTok{)}
\NormalTok{  pre }\OtherTok{\textless{}{-}} \FunctionTok{rnorm}\NormalTok{(n, }\AttributeTok{mean =} \DecValTok{70}\NormalTok{, }\AttributeTok{sd =}\NormalTok{ sd)}
\NormalTok{  post }\OtherTok{\textless{}{-}} \DecValTok{72} \SpecialCharTok{+}
\NormalTok{    rho }\SpecialCharTok{*}\NormalTok{ (pre }\SpecialCharTok{{-}} \DecValTok{70}\NormalTok{) }\SpecialCharTok{+}
\NormalTok{    effect }\SpecialCharTok{*}\NormalTok{ treat }\SpecialCharTok{+}
    \FunctionTok{rnorm}\NormalTok{(n, }\AttributeTok{mean =} \DecValTok{0}\NormalTok{, }\AttributeTok{sd =}\NormalTok{ sd }\SpecialCharTok{*} \FunctionTok{sqrt}\NormalTok{(}\DecValTok{1} \SpecialCharTok{{-}}\NormalTok{ rho}\SpecialCharTok{\^{}}\DecValTok{2}\NormalTok{))}
  \FunctionTok{data.frame}\NormalTok{(treat, pre, post)}
\NormalTok{\}}

\NormalTok{analyze\_ancova }\OtherTok{\textless{}{-}} \ControlFlowTok{function}\NormalTok{(dat) \{}
  \FunctionTok{summary}\NormalTok{(}\FunctionTok{lm}\NormalTok{(post }\SpecialCharTok{\textasciitilde{}}\NormalTok{ treat }\SpecialCharTok{+}\NormalTok{ pre, }\AttributeTok{data =}\NormalTok{ dat))}\SpecialCharTok{$}\NormalTok{coefficients[}\StringTok{"treat"}\NormalTok{, ]}
\NormalTok{\}}
\end{Highlighting}
\end{Shaded}
\end{frame}

\begin{frame}[fragile]{What if we make this more complicated,
arbitrarily so?}
\phantomsection\label{what-if-we-make-this-more-complicated-arbitrarily-so-2}
\begin{Shaded}
\begin{Highlighting}[]
\FunctionTok{set.seed}\NormalTok{(}\DecValTok{2026}\NormalTok{)}
\FunctionTok{simulate\_power}\NormalTok{(}\DecValTok{2000}\NormalTok{, gen\_pretest, analyze\_lm, success\_sig)}
\end{Highlighting}
\end{Shaded}

\begin{verbatim}
[1] 0.807
\end{verbatim}

\begin{Shaded}
\begin{Highlighting}[]
\FunctionTok{simulate\_power}\NormalTok{(}\DecValTok{2000}\NormalTok{, gen\_pretest, analyze\_ancova, success\_sig)}
\end{Highlighting}
\end{Shaded}

\begin{verbatim}
[1] 0.9415
\end{verbatim}

\vskip10pt

\begin{itemize}
\tightlist
\item
  We can see that our power jumps from about 80\% to about \textbf{94\%}
\item
  Power is a property of the
  \textcolor{red}{design \textit{and} the analysis}, not just of \(N\)
\item
  If we were really strapped for students to enroll, this could be
  really helpful to know!
\end{itemize}

\small

(To be fair: this case still has a closed form (\(\sigma\) just becomes
\(\sigma\sqrt{1-\rho^2}\)) but we didn't need to know that!)
\end{frame}

\begin{frame}{Let's make it even \textbf{more} complicated!}
\phantomsection\label{lets-make-it-even-more-complicated}
Consider the classic problem of dropout. Let's make the following
assumptions about our study (maybe we base this on some sort of past
experience running such a study):

\begin{itemize}
\tightlist
\item
  Students who are doing \emph{badly} are more likely to stop showing up
\item
  The \textbf{control group loses more students} (here, about 23\%
  vs.~11\%)
\end{itemize}

\vskip20pt

\begin{center}
You may ask: \textit{Is there a formula for the power of this study?} \\
\vskip5pt
\textbf{Not really.} Dropout depends on an \textit{unrealized outcome}, so the \textcolor{red}{sample size}, the \textcolor{red}{effect we estimate}, \textit{and} the \textcolor{red}{noise} all change at once!! What a pain!!
\end{center}
\end{frame}

\begin{frame}[fragile]{Let's make it even \textbf{more} complicated!}
\phantomsection\label{lets-make-it-even-more-complicated-1}
But we don't need a formula! We can use our computers!! We just need to
describe how the data are generated:

\begin{Shaded}
\begin{Highlighting}[]
\NormalTok{gen\_dropout }\OtherTok{\textless{}{-}} \ControlFlowTok{function}\NormalTok{(}\AttributeTok{n =} \DecValTok{350}\NormalTok{, }\AttributeTok{effect =} \DecValTok{3}\NormalTok{, }\AttributeTok{sd =} \DecValTok{10}\NormalTok{, }\AttributeTok{rho =} \FloatTok{0.6}\NormalTok{) \{}
\NormalTok{    dat }\OtherTok{\textless{}{-}} \FunctionTok{gen\_pretest}\NormalTok{(n, effect, sd, rho)}
\NormalTok{    p\_drop }\OtherTok{\textless{}{-}} \FunctionTok{plogis}\NormalTok{(}\SpecialCharTok{{-}}\FloatTok{2.2} \SpecialCharTok{+} \FloatTok{0.7} \SpecialCharTok{*}\NormalTok{ (dat}\SpecialCharTok{$}\NormalTok{treat }\SpecialCharTok{==} \DecValTok{0}\NormalTok{) }\SpecialCharTok{{-}} \FloatTok{0.12} \SpecialCharTok{*}\NormalTok{ (dat}\SpecialCharTok{$}\NormalTok{post }\SpecialCharTok{{-}} \DecValTok{72}\NormalTok{))}
\NormalTok{    dat[}\FunctionTok{runif}\NormalTok{(n) }\SpecialCharTok{\textgreater{}}\NormalTok{ p\_drop, ] }\CommentTok{\# we only observe the students who stayed}
\NormalTok{\}}

\FunctionTok{nrow}\NormalTok{(}\FunctionTok{gen\_dropout}\NormalTok{()) }\CommentTok{\# out of the 350 we enrolled}
\end{Highlighting}
\end{Shaded}

\begin{verbatim}
[1] 289
\end{verbatim}
\end{frame}

\begin{frame}[fragile]{Let's make it even \textbf{more} complicated!}
\phantomsection\label{lets-make-it-even-more-complicated-2}
\begin{Shaded}
\begin{Highlighting}[]
\FunctionTok{set.seed}\NormalTok{(}\DecValTok{2026}\NormalTok{)}
\FunctionTok{simulate\_power}\NormalTok{(}\DecValTok{2000}\NormalTok{, gen\_pretest, analyze\_ancova, success\_sig)}
\end{Highlighting}
\end{Shaded}

\begin{verbatim}
[1] 0.9425
\end{verbatim}

\begin{Shaded}
\begin{Highlighting}[]
\FunctionTok{simulate\_power}\NormalTok{(}\DecValTok{2000}\NormalTok{, gen\_dropout, analyze\_ancova, success\_sig)}
\end{Highlighting}
\end{Shaded}

\begin{verbatim}
[1] 0.6815
\end{verbatim}

\vskip10pt

\begin{itemize}
\tightlist
\item
  Remember, dropout doesn't just shrink our sample!
\end{itemize}
\end{frame}

\begin{frame}[fragile]{Let's make it even \textbf{more} complicated!}
\phantomsection\label{lets-make-it-even-more-complicated-3}
Because we simulated the data, we \emph{know} the true effect is 3. So
we can check:

\begin{Shaded}
\begin{Highlighting}[]
\FunctionTok{set.seed}\NormalTok{(}\DecValTok{2026}\NormalTok{)}
\NormalTok{estimates }\OtherTok{\textless{}{-}} \FunctionTok{replicate}\NormalTok{(}\DecValTok{2000}\NormalTok{, }\FunctionTok{analyze\_ancova}\NormalTok{(}\FunctionTok{gen\_dropout}\NormalTok{())[[}\StringTok{"Estimate"}\NormalTok{]])}
\FunctionTok{mean}\NormalTok{(estimates)}
\end{Highlighting}
\end{Shaded}

\begin{verbatim}
[1] 2.229586
\end{verbatim}

\vskip10pt

\begin{itemize}
\tightlist
\item
  Since students doing poorly \emph{leave}, this biases our results for
  the control group up
\item
  We appear to \textbf{underestimate} the effect
\end{itemize}

\begin{center}
Unfortuantely, there's no formula for this!!
\end{center}
\end{frame}

\begin{frame}[fragile]{Let's make it even \textbf{more} complicated!}
\phantomsection\label{lets-make-it-even-more-complicated-4}
Remember our relaxed definition:
\textbf{the probability that we \textcolor{red}{recover} a real effect, should one \textcolor{red}{exist}}.

\begin{itemize}
\tightlist
\item
  ``Recover'' doesn't have to mean a p-value.
\end{itemize}

\begin{Shaded}
\begin{Highlighting}[]
\NormalTok{success\_close }\OtherTok{\textless{}{-}} \ControlFlowTok{function}\NormalTok{(est) }\FunctionTok{abs}\NormalTok{(est[[}\StringTok{"Estimate"}\NormalTok{]] }\SpecialCharTok{{-}} \DecValTok{3}\NormalTok{) }\SpecialCharTok{\textless{}} \FloatTok{1.5}

\FunctionTok{set.seed}\NormalTok{(}\DecValTok{2026}\NormalTok{)}
\FunctionTok{simulate\_power}\NormalTok{(}\DecValTok{2000}\NormalTok{, gen\_pretest, analyze\_ancova, success\_close) }\CommentTok{\# no dropout}
\end{Highlighting}
\end{Shaded}

\begin{verbatim}
[1] 0.9185
\end{verbatim}

\begin{Shaded}
\begin{Highlighting}[]
\FunctionTok{simulate\_power}\NormalTok{(}\DecValTok{2000}\NormalTok{, gen\_dropout, analyze\_ancova, success\_close) }\CommentTok{\# dropout}
\end{Highlighting}
\end{Shaded}

\begin{verbatim}
[1] 0.786
\end{verbatim}
\end{frame}

\begin{frame}{Let's make it even \textbf{more} complicated!}
\phantomsection\label{lets-make-it-even-more-complicated-5}
\large

There are other options for recovery! You could ask when:

\begin{itemize}
\tightlist
\item
  the CI excludes zero
\item
  the CI is narrower than some width
\item
  the sign is right
\item
  a posterior probability exceeds 0.9 or 0.95
\item
  and much more!
\end{itemize}
\end{frame}

\begin{frame}[fragile]{So, how many students do we need?!}
\phantomsection\label{so-how-many-students-do-we-need}
Just like before, we want to find the \(N\) that gets us to 80\%. Now,
instead of inverting a formula, we try a grid of sample sizes:

\begin{Shaded}
\begin{Highlighting}[]
\FunctionTok{set.seed}\NormalTok{(}\DecValTok{2026}\NormalTok{)}
\NormalTok{ns }\OtherTok{\textless{}{-}} \FunctionTok{seq}\NormalTok{(}\DecValTok{150}\NormalTok{, }\DecValTok{500}\NormalTok{, }\AttributeTok{by =} \DecValTok{50}\NormalTok{)}
\NormalTok{powers }\OtherTok{\textless{}{-}} \FunctionTok{sapply}\NormalTok{(ns, \textbackslash{}(n) \{}
  \FunctionTok{simulate\_power}\NormalTok{(}
    \DecValTok{1000}\NormalTok{,}
    \AttributeTok{generate =}\NormalTok{ \textbackslash{}() }\FunctionTok{gen\_dropout}\NormalTok{(}\AttributeTok{n =}\NormalTok{ n),}
    \AttributeTok{analyze =}\NormalTok{ analyze\_ancova,}
    \AttributeTok{success =}\NormalTok{ success\_sig}
\NormalTok{  )}
\NormalTok{\})}
\end{Highlighting}
\end{Shaded}
\end{frame}

\begin{frame}{So, how many students do we need?}
\phantomsection\label{so-how-many-students-do-we-need-1}
\includegraphics{statistical-power-combined_files/figure-beamer/unnamed-chunk-14-1.pdf}

\begin{center}
\small With dropout, we need to enrol \textbf{around 450 students}. Alas, more students only buys \textit{power} because the estimate is \textcolor{red}{still biased}.
\end{center}
\end{frame}

\begin{frame}{Your turn!}
\phantomsection\label{your-turn}
We provided you with a template file to use for htis part,
\texttt{power-sim-template.R}. Use this and fill in every
\texttt{\_\_\_}. Use your own study, or, if nothing comes to mind, the
tutoring study from earlier is good practice too!

Ask yourself:

\begin{enumerate}
\tightlist
\item
  What are my baseline, effect, and noise values, and where did I get
  those numbers?
\item
  \texttt{generate()} -- Does this dataset look like something my study
  could produce?
\item
  \texttt{analyze()} -- What would I \textit{actually} run on the real
  data?
\item
  \texttt{success()} -- What counts as ``recovering'' the effect?
\item
  Estimate power -- With \textit{no} effect, do I get about 5\%? A
  significance-based \texttt{success()} should succeed about 5\% of the
  time
\item
  Perform a sweep over \(N\) -- Where does my curve cross 80\%?
\item
  Now, \textcolor{statsecondary}{\textbf{add a complication}} -- Maybe
  unequal groups, a covariate, a skewed outcome, with these how much
  power do I lose?
\end{enumerate}
\end{frame}



\end{document}
